\section{Conclusion}
\label{sec:conclusion}

We began by asking whether better AI makes us intellectually lazier. Our answer is: for returning WildChat-1M users in 2023--2024, empirically, no --- prompt complexity grows significantly ($\tau = +0.744$, $p < 0.001$), not declines. But the right comparison group has not yet been assembled, the measurement infrastructure to build that comparison fully is not yet publicly available, and --- importantly --- the causal chain from AI quality improvement to behavioral change has not been verified. The result refutes the BAA directional prediction for prompt length in this cohort; it does not close the BAA empirical question.

We contributed four things: the first large-scale empirical test of the BAA disengagement directional prediction (to our knowledge, pending citation verification; negative result); a validated behavioral trend measurement pipeline (Hamed-Rao Mann-Kendall on returning-user cohorts from WildChat-1M); documentation of a binding LMSYS data access constraint; and characterization of the selection bias problem with a specific proposed remedy.

The most critical next steps are a two-group comparison (returning vs.\ non-returning users, difference-in-differences), LMSYS primary access for vote entropy, inclusion of an AI quality covariate per interaction bin, and implicit correction proxy development. The question of whether AI improvement changes human behavior --- and in what direction --- is answerable from existing interaction data. We have established the methodological foundation and documented the empirical boundary; resolving the ambiguity requires better data access, a more careful cohort design, and per-interaction AI quality measures not currently available in public datasets.
